% ******************************* PhD Thesis Template **************************
% Please have a look at the README.md file for info on how to use the template

\PassOptionsToPackage{style=apa, backend=biber}{biblatex}

\documentclass[a4paper, 12pt, print, index, oneside, times, custombib]{PhDThesisPSnPDF}


%\defbibenvironment{bibliography}
%{\list
	%	{}
	%	{\setlength{\leftmargin}{\bibhang}%
		%		\setlength{\itemindent}{-\leftmargin}%
		%		\setlength{\itemsep}{\bibitemsep}%
		%		\setlength{\parsep}{\bibparsep}}}
%{\endlist}
%{\item}
\usepackage[style=apa, backend=biber]{biblatex}
\addbibresource{Referencias.bib}


% Definición de encabezamientos y pié de página

\usepackage{fancyhdr}
\pagestyle{fancy}
\fancyhf{}

\lhead[\smaller Nombre]{\smaller Nombre}

\fancyhead[R]{\smaller Título}

\fancyfoot[CE, CO]{\thepage}

\usepackage[ruled,vlined]{algorithm2e}
\usepackage[T1]{fontenc}
\usepackage{lmodern}
\usepackage{booktabs}
\usepackage{xcolor}
\usepackage[none]{hyphenat}
\usepackage{longtable}
\usepackage[utf8]{inputenc}
\usepackage{array}       % mejor control de columnas
\usepackage{geometry}    % márgenes
\usepackage{float}
\usepackage{subcaption} %recortar palabras y que entren en los margenes
\usepackage[spanish]{babel}
\usepackage{microtype}
\usepackage{hyphenat}
\sloppy
\emergencystretch=2em
\usepackage{needspace}


%********************************
\usepackage{babel}
\addto\captionsspanish{
  \renewcommand{\listtablename}{Índice de tablas}
  \renewcommand{\tablename}{Tabla}
}
\renewcaptionname{spanish}{\tablename}{Tabla}

\title{Título}

%% Subtitle (Optional)

%\subtitle{Subtítulo}
%\else
%\subtitle{Using the ULL template}
%\fi

%% The full name of the author
\author{Nombre y Apellido}
%\date{}

%% Department (eg. Department of Engineering, Maths, Physics)

\dept{Departamento de Sistemas Digitales y Datos}


%% University and Crest

\university{Instituto Tecnológico de Buenos Aires}

\crest{\includegraphics[width=0.5\textwidth]{derf__64049_1952016_LOGOITBA.jpg}}





%% Supervisor (optional)
%% for multiple supervisors, append each supervisor with the \newline command
\supervisor{Tutor/a}




\degreetitle{Master en Ciencias de Datos}




\ifdefineAbstract
\includeonly{Abstract/abstract} 



\fi
% ***************************** Chapter Mode ***********************************
% The chapter mode allows user to only print particular chapters with references
% Title, Contents, Frontmatter are disabled by default
% Useful option to review a particular chapter or to send it to supervisior.
% To use choose `chapter' option in the document class

\ifdefineChapter
%\includeonly{Chapter3/chapter3}
\fi

% ******************************** Front Matter ********************************

\renewcommand{\contentsname}{\'Indice General}
\renewcommand{\listfigurename}{\'Indice de Figuras}
\renewcommand{\listtablename}{\'Indice de Tablas}
\renewcommand{\chaptername}{Capítulo}
\renewcommand{\bibname}{Bibliografía}
\renewcommand{\figurename}{Figura}
\renewcommand{\tablename}{Tabla}

%%%%%%%%%%%%%%%%%%%%%%%%%%%%%%%%%%%%%%%%%%%%%%%%%%%%%%%%%%%%%%%%%%%

\begin{document}

\frontmatter

\maketitle

%	\include{dedication}
%	\include{Declaration/declaration}
%	\include{Acknowledgement/acknowledgement}
%	\include{Abstract/abstract}

% *********************** Adding TOC and List of Figures ***********************

\tableofcontents

\listoffigures

\listoftables

% \printnomenclature[space] space can be set as 2em between symbol and description
%\printnomenclature[3em]

\printnomenclature

% ******************************** Main Matter *********************************
\mainmatter
%	\include{Abstract/abstract}
%	\include{Chapters/chapter1}
%	\include{Chapters/chapter2}
%        \include{Chapters/chapter3}
%        \include{Chapters/chapter4}
%\include{Chapters/chapter4}
%\include{Chapters/chapter5}
%\include{Chapters/chapter6}
%\include{Chapters/chapter7}

%%%%%%%%%%%%%%%%%%%%%%%%%%%%%%%%%%%%%%%%%%%%%%%%%%%%%%%%%%%%%%%%%%%%%%%%%%%%%%%%%
\include{dedication}

	

%%%%%%%%%%%%%%%%%%%%%%%%%%%%%%%%%%%%%%%%%%%%%%%%%%%%%%%%%%%%%%%%%%%%%%%%%%%%%%%%%


\chapter{Introducción}
\label{sec:Introducción}

Este en un ejemplo de cómo citar una referencia~\parencite{MachineLearningandDeepLearning}.

\textcite{MachineLearningandDeepLearning} proponen...


\textcolor{red}{Incluir la estructura del trabajo describiendo brevemente lo que se va a desarrollar en cada capítulo.}

%%%%%%%%%%%%%%%%%%%%%%%%%%%%%%%%%%%%%%%%%%%%%%%%%%%%%%%%%%%%%%%%%%%%%%%%%%%%%%%%%
\section{Planteamiento del problema} 
\label{sec:Planteamiento}



%%%%%%%%%%%%%%%%%%%%%%%%%%%%%%%%%%%%%%%%%%%%%%%%%%%%%%%%%%%%%%%%%%%%%%%%%%%%%%%%%
\section{Justificación del estudio}
\label{sec:Justificación}


%%%%%%%%%%%%%%%%%%%%%%%%%%%%%%%%%%%%%%%%%%%%%%%%%%%%%%%%%%%%%%%%%%%%%%%%%%%%%%%%%
\section{Hipótesis de trabajo}
\label{sec:Hipótesis}

%%%%%%%%%%%%%%%%%%%%%%%%%%%%%%%%%%%%%%%%%%%%%%%%%%%%%%%%%%%%%%%%%%%%%%%%%%%%%%%%%
\section{Objetivos}
\subsection{Objetivo general}
\label{sec:Objetivo}



\subsection{Objetivos específicos}
\label{sec:Objetivos}


%%%%%%%%%%%%%%%%%%%%%%%%%%%%%%%%%%%%%%%%%%%%%%%%%%%%%%%%%%%%%%%%%%%%%%%%%%%%%%%%%
\section{Estado del arte}
\label{sec:SOTA}


%%%%%%%%%%%%%%%%%%%%%%%%%%%%%%%%%%%%%%%%%%%%%%%%%%%%%%%%%%%%%%%%%%%%%%%%%%%%%%%%%
\chapter{Metodología}
\label{sec:Metodología}

\section{Conjunto de datos}
\label{sec:Datos}

%%%%%%%%%%%%%%%%%%%%%%%%%%%%%%%%%%%%%%%%%%%%%%%%%%%%%%%%%%%%%%%%%%%%%%%%%%%%%%%%
\section{Modelos}
\label{sec:Modelos}

%%%%%%%%%%%%%%%%%%%%%%%%%%%%%%%%%%%%%%%%%%%%%%%%%%%%%%%%%%%%%%%%%%%%%%%%%%%%%%%%

\section{Métricas de evaluación}
\label{sec:Métricas}


%%%%%%%%%%%%%%%%%%%%%%%%%%%%%%%%%%%%%%%%%%%%%%%%%%%%%%%%%%%%%%%%%%%%%%%%%%%%%%%%%

\chapter{Resultados}
\label{sec:Resultados}

%%%%%%%%%%%%%%%%%%%%%%%%%%%%%%%%%%%%%%%%%%%%%%%%%%%%%%%%%%%%%%%%%%%%%%%%%%%%%%%%%

\chapter{Discusión}
\label{sec:Discusión}

%%%%%%%%%%%%%%%%%%%%%%%%%%%%%%%%%%%%%%%%%%%%%%%%%%%%%%%%%%%%%%%%%%%%%%%%%%%%%%%%%

\chapter{Conclusiones}
\label{sec:Conclusiones}

%%%%%%%%%%%%%%%%%%%%%%%%%%%%%%%%%%%%%%%%%%%%%%%%%%%%%%%%%%%%%%%%%%%%%%%%%%%%%%%%%

% ********************************** Back Matter *******************************
% Backmatter should be commented out, if you are using appendices after References
%\backmatter

% ********************************** Bibliography ******************************
%\begin{spacing}{0.9}
	
	% To use the conventional natbib style referencing
	% Bibliography style previews: http://nodonn.tipido.net/bibstyle.php
	% Reference styles: http://sites.stat.psu.edu/~surajit/present/bib.htm
	
	%\bibliographystyle{apalike}
	%\bibliographystyle{unsrt} % Use for unsorted references  
	%	\bibliographystyle{plainnat} % use this to have URLs listed in References
	%	\cleardoublepage
	%	\bibliography{References} % Path to your References.bib file
	
	
	% If you would like to use BibLaTeX for your references, pass `custombib' as
	% an option in the document class. The location of 'reference.bib' should be
	% specified in the preamble.tex file in the custombib section.
	% Comment out the lines related to natbib above and uncomment the following line.
	
	%\printbibliography[heading=bibintoc, title={References}]
	
	
%\end{spacing}

\printbibliography[heading=bibintoc, title={Bibliogafía}]

% ********************************** Appendices ********************************

\begin{appendices} % Using appendices environment for more functunality
	
	%\include{appendix1}
%	\include{Appendix2/appendix2}
	
\end{appendices}

% *************************************** Index ********************************
\printthesisindex % If index is present
%	
%



\end{document}
